% Azure section (App Service, SDK/CLI, AI Foundry)

% 1 – Azure App Service
\QuestionSlide[\CategoryBadge[AzureColor!20]{Azure}]{* What is Azure App Service?}

\begin{frame}[fragile]
  \frametitle{%
    \begin{tikzpicture}[remember picture, overlay]
      \node[anchor=north east, xshift=-0.4cm, yshift=-0.4cm, text=black] at (current page.north east) {
        \CategoryBadge[AzureColor!20]{Azure}
      };
    \end{tikzpicture}
    Answer \theqcounter: What is Azure App Service?%
  }

  {\footnotesize
  \textbf{Azure App Service} is a fully managed PaaS for hosting web apps, REST APIs, and WebJobs. Supports .NET, Node.js, Python, Java, and containers. Key features: \textbf{deployment slots} for staged rollouts, \textbf{autoscaling}, \textbf{managed TLS}, VNet integration, and built-in monitoring via Application Insights.
  }

  \begin{minted}[fontsize=\scriptsize]{bash}
# Deploy from GitHub Actions / Azure CLI
az webapp create \
    --resource-group myRG \
    --plan myPlan \
    --name myapi-prod \
    --runtime "DOTNET:9"

az webapp deployment source config-zip \
    --resource-group myRG --name myapi-prod \
    --src publish.zip
  \end{minted}
\end{frame}

% 2 – Managed Identity
\QuestionSlide[\CategoryBadge[AzureColor!20]{Azure}]{** What is Azure Managed Identity?}

\begin{frame}[fragile]
  \frametitle{%
    \begin{tikzpicture}[remember picture, overlay]
      \node[anchor=north east, xshift=-0.4cm, yshift=-0.4cm, text=black] at (current page.north east) {
        \CategoryBadge[AzureColor!20]{Azure}
      };
    \end{tikzpicture}
    Answer \theqcounter: What is Managed Identity?%
  }

  {\footnotesize
  \textbf{Managed Identity} gives an Azure resource (App Service, AKS Pod, VM) an AAD identity without storing credentials. The Azure platform issues and rotates tokens automatically. Use \texttt{DefaultAzureCredential} in the Azure SDK — it tries Managed Identity in production and developer credentials locally, requiring zero config changes between environments.
  }

  \begin{minted}[fontsize=\scriptsize]{csharp}
// Same code works locally (developer login) and in Azure (Managed Identity)
var credential = new DefaultAzureCredential();

var secretClient = new SecretClient(
    new Uri("https://myvault.vault.azure.net/"),
    credential);

KeyVaultSecret secret = await secretClient.GetSecretAsync("DbPassword");
  \end{minted}
\end{frame}

% 3 – Azure Key Vault
\QuestionSlide[\CategoryBadge[AzureColor!20]{Azure}]{* What is Azure Key Vault?}

\begin{frame}[fragile]
  \frametitle{%
    \begin{tikzpicture}[remember picture, overlay]
      \node[anchor=north east, xshift=-0.4cm, yshift=-0.4cm, text=black] at (current page.north east) {
        \CategoryBadge[AzureColor!20]{Azure}
      };
    \end{tikzpicture}
    Answer \theqcounter: What is Azure Key Vault?%
  }

  {\footnotesize
  \textbf{Azure Key Vault} is a managed HSM-backed service for storing \textbf{secrets} (connection strings, API keys), \textbf{certificates}, and \textbf{cryptographic keys}. Access is controlled via AAD roles (RBAC) or access policies. Integrate with ASP.NET Core \texttt{IConfiguration} to surface secrets as configuration values without code changes.
  }

  \begin{minted}[fontsize=\scriptsize]{csharp}
builder.Configuration.AddAzureKeyVault(
    new Uri("https://myapp-kv.vault.azure.net/"),
    new DefaultAzureCredential());

// Secret "ConnectionStrings--DefaultConnection" is now available as:
var connStr = builder.Configuration.GetConnectionString("DefaultConnection");
  \end{minted}
\end{frame}

% 4 – Azure AI Foundry
\QuestionSlide[\CategoryBadge[AzureColor!20]{Azure}]{** What is Azure AI Foundry?}

\begin{frame}[fragile]
  \frametitle{%
    \begin{tikzpicture}[remember picture, overlay]
      \node[anchor=north east, xshift=-0.4cm, yshift=-0.4cm, text=black] at (current page.north east) {
        \CategoryBadge[AzureColor!20]{Azure}
      };
    \end{tikzpicture}
    Answer \theqcounter: What is Azure AI Foundry?%
  }

  {\footnotesize
  \textbf{Azure AI Foundry} (formerly Azure AI Studio) is Microsoft's unified platform for building, evaluating, and deploying AI applications. It provides: a \textbf{model catalogue} (OpenAI GPT-4o, Phi, Llama), \textbf{Prompt Flow} for LLM pipeline authoring, \textbf{evaluation} dashboards for quality/safety, and \textbf{deployment} to managed endpoints behind Azure OpenAI APIs.
  }
\end{frame}

% 5 – Azure AI Foundry Prompt Flow
\QuestionSlide[\CategoryBadge[AzureColor!20]{Azure}]{** What is Azure AI Foundry Prompt Flow?}

\begin{frame}[fragile]
  \frametitle{%
    \begin{tikzpicture}[remember picture, overlay]
      \node[anchor=north east, xshift=-0.4cm, yshift=-0.4cm, text=black] at (current page.north east) {
        \CategoryBadge[AzureColor!20]{Azure}
      };
    \end{tikzpicture}
    Answer \theqcounter: What is Prompt Flow?%
  }

  {\footnotesize
  \textbf{Prompt Flow} is a low-code / code-first tool for building LLM pipelines as directed acyclic graphs (DAGs) of \textbf{nodes}: LLM calls, Python code, conditional branches, and vector index lookups. Flows can be tested, evaluated against ground-truth datasets, and deployed as streaming HTTP endpoints. SDK available for .NET and Python.
  }

  \begin{minted}[fontsize=\scriptsize]{python}
# Prompt Flow SDK – run a deployed flow from .NET via HTTP
# or call directly in Python
from promptflow import PFClient
client = PFClient()
result = client.run(
    flow="./my_rag_flow",
    data="./eval_data.jsonl",
    column_mapping={"question": "${data.question}"}
)
  \end{minted}
\end{frame}

% 6 – Azure Service Bus
\QuestionSlide[\CategoryBadge[AzureColor!20]{Azure}]{** What is Azure Service Bus and when do you use it?}

\begin{frame}[fragile]
  \frametitle{%
    \begin{tikzpicture}[remember picture, overlay]
      \node[anchor=north east, xshift=-0.4cm, yshift=-0.4cm, text=black] at (current page.north east) {
        \CategoryBadge[AzureColor!20]{Azure}
      };
    \end{tikzpicture}
    Answer \theqcounter: What is Azure Service Bus?%
  }

  {\footnotesize
  \textbf{Azure Service Bus} is a fully managed enterprise message broker. Use \textbf{Queues} for point-to-point async communication; \textbf{Topics + Subscriptions} for pub/sub fan-out. Key features: message sessions, dead-letter queue, scheduled delivery, and exactly-once delivery with duplicate detection. Ideal for decoupling microservices and smoothing traffic spikes.
  }

  \begin{minted}[fontsize=\scriptsize]{csharp}
await using var client = new ServiceBusClient(connectionString);
ServiceBusSender sender = client.CreateSender("orders");
await sender.SendMessageAsync(new ServiceBusMessage(
    BinaryData.FromObjectAsJson(new { OrderId = 42 })));

// Consumer
ServiceBusProcessor proc = client.CreateProcessor("orders");
proc.ProcessMessageAsync += async args =>
    { await HandleOrder(args.Message); await args.CompleteMessageAsync(args.Message); };
await proc.StartProcessingAsync();
  \end{minted}
\end{frame}

% 7 – Azure CLI essential commands
\QuestionSlide[\CategoryBadge[AzureColor!20]{Azure}]{* What are essential Azure CLI patterns for .NET developers?}

\begin{frame}[fragile]
  \frametitle{%
    \begin{tikzpicture}[remember picture, overlay]
      \node[anchor=north east, xshift=-0.4cm, yshift=-0.4cm, text=black] at (current page.north east) {
        \CategoryBadge[AzureColor!20]{Azure}
      };
    \end{tikzpicture}
    Answer \theqcounter: Essential Azure CLI patterns?%
  }

  {\footnotesize
  The Azure CLI (\texttt{az}) automates provisioning and deployment from scripts and pipelines. Authenticate via \texttt{az login} or a Service Principal in CI. Use \texttt{--output table/json/tsv} for scripting.
  }

  \begin{minted}[fontsize=\scriptsize]{bash}
az login                                         # interactive
az account set --subscription "My Subscription"

# Provision an App Service
az group create -n myRG -l westeurope
az appservice plan create -n myPlan -g myRG --sku B2
az webapp create -n myapi -g myRG -p myPlan --runtime "DOTNET:9"

# Set app settings / connection strings
az webapp config appsettings set -n myapi -g myRG \
    --settings ASPNETCORE_ENVIRONMENT=Production

# Tail live logs
az webapp log tail -n myapi -g myRG
  \end{minted}
\end{frame}

% 8 – Azure Monitor & Application Insights
\QuestionSlide[\CategoryBadge[AzureColor!20]{Azure}]{** What is Application Insights and how do you add it to an ASP.NET Core app?}

\begin{frame}[fragile]
  \frametitle{%
    \begin{tikzpicture}[remember picture, overlay]
      \node[anchor=north east, xshift=-0.4cm, yshift=-0.4cm, text=black] at (current page.north east) {
        \CategoryBadge[AzureColor!20]{Azure}
      };
    \end{tikzpicture}
    Answer \theqcounter: What is Application Insights?%
  }

  {\footnotesize
  \textbf{Application Insights} (part of Azure Monitor) collects traces, metrics, exceptions, dependencies, and live metrics from your app. \texttt{AddApplicationInsightsTelemetry()} wires it into ASP.NET Core with zero code changes. It integrates with OpenTelemetry exporters so the same instrumentation works locally and in Azure.
  }

  \begin{minted}[fontsize=\scriptsize]{csharp}
// Program.cs
builder.Services.AddApplicationInsightsTelemetry(
    builder.Configuration["APPLICATIONINSIGHTS_CONNECTION_STRING"]);

// Alternatively, using OpenTelemetry exporter:
builder.Services.AddOpenTelemetry()
    .WithTracing(t => t.AddAspNetCoreInstrumentation()
                       .AddAzureMonitorTraceExporter());
  \end{minted}
\end{frame}
